\chapter{Accounting And Product Semantics}

\section{Logical Size First}

The v1 metric is logical size. This is the only size metric with enough current
evidence for the narrow raw parser contract. Directory aggregates use
unique-inode logical totals within the aggregate root so hard links do not
double-count inside one aggregate.

Source precedence for logical size is fixed by
\path{docs/research/sources/SR-009-compression-logical-size-precedence.md} and
the wider, scoped table in
\path{docs/research/sources/SR-017-logical-size-source-precedence.md}. The
current accounting log is
\path{docs/research/RL-07-size-and-space-accounting.md}.

\section{Deferred Metrics}

The following metrics remain research topics:

\begin{itemize}
  \item allocated size
  \item exclusive size
  \item shared size
  \item compression storage savings
  \item snapshot-retained bytes
\end{itemize}

\path{docs/research/experiments/EX-09-accounting-probe/README.md} defines the
next accounting probe. It explicitly avoids treating \texttt{du},
\texttt{st\_blocks}, raw extents, and Finder-like values as interchangeable
oracles.

\section{Boot-Root Semantics}

Raw single-volume namespace and user-visible macOS root namespace are distinct.
Modern macOS may involve System/Data volume groups, firmlinks, snapshots, and
signed system volume behavior. The manual should never imply that a raw
single-volume walk is the same as Finder-visible \texttt{/}.

The product-semantics log is
\path{docs/research/RL-11-snapshots-volume-groups-and-firmlinks.md} and the
boundary source review is
\path{docs/research/sources/SR-010-snapshots-volume-groups-firmlinks.md}.

\section{Long-Range Roadmap}

The general WizTree-for-any-Mac product path is tracked in
\path{docs/research/plans/general-wiztree-for-any-mac-roadmap.md}. It defines
staged gates from the narrow Rust full scan to a hybrid consumer product:

\begin{description}
  \item[Gate 0] QA discipline (artifacts, oracles, fail-closed checks).
  \item[Gate 1] Native narrow Rust full scan
  (\texttt{raw\_single\_volume\_logical\_size}).
  \item[Gate 2] Local single-volume backend with raw fast path and supported-API
  fallback.
  \item[Gate 3] Finder-visible macOS namespace mode
  (\texttt{os\_visible\_namespace}).
  \item[Gate 4] Encryption and live runtime support.
  \item[Gate 5] Size semantics beyond logical size.
  \item[Gate 6] Incremental scanning and persistent cache.
  \item[Gate 7] Performance engineering.
  \item[Gate 8] Product packaging and UX.
\end{description}

A parallel release shape (R0 research harness, R1 narrow Rust MWP, R2 hybrid
single-volume scanner, R3 macOS visible namespace scanner, R4 fast repeat-scan
product) keeps the audience and the support boundary explicit at every release.
The roadmap's blocker register names \texttt{EX-14}, the SR-015 xfield replay,
the Rust body-field dump, the compressed logical-size fixture, and the
live-startup support questions as the current open work.
